% Portada de C5: el ciclo de desarrollo con las pruebas destacadas.
\documentclass[tikz,border=0pt]{standalone}
\usepackage{resonancia}
\begin{document}
\begin{tikzpicture}
  \clip[rounded corners=6pt] (0,0) rectangle (12,3.6);
  \shade[top color=rGris!22, bottom color=white] (0,0) rectangle (12,3.6);
  \foreach \i/\t/\c in {0/explorar/rAzul, 1/planificar/rAzul, 2/implementar/rNaranja, 3/probar/rVerde, 4/revisar/rVioleta, 5/documentar/rGris, 6/commit/rGris} {
    \node[caja, draw=\c, fill=white, minimum width=1.45cm, minimum height=0.7cm, inner sep=2pt, font=\sffamily\scriptsize, text=\c] (b\i) at ({0.95+1.65*\i},1.8) {\t};
  }
  \draw[rVerde, line width=2pt, rounded corners=2pt] ($(b3.south west)+(-0.05,-0.05)$) rectangle ($(b3.north east)+(0.05,0.05)$);
  \foreach \i/\j in {0/1,1/2,2/3,3/4,4/5,5/6} { \draw[trazo, color=rGris, -{Stealth}] (b\i) -- (b\j); }
  \draw[trazo, color=rVerde, -{Stealth}] (b3.south) .. controls (6.6,0.6) and (3.4,0.6) .. (b2.south);
  \node[etiqueta, color=rVerde] at (4.9,0.5) {si falla, vuelve};
  \node[etiqueta, color=rGris] at (6,3.1) {tú decides qué es correcto};
\end{tikzpicture}
\end{document}
